\section{Optimizing the scale profile}
\label{humanprofile:section}

This section supplies the real-variable part of the proof. A profile records
how the masses of the regularized pin measure change with scale. The analytic
argument needs a decreasing sequence of scales with small total cost; the
purpose here is to construct that sequence and evaluate the resulting
dimension condition explicitly.

\subsection{The cost to be bounded}

Let \(g:[0,N]\to\mathbb R\) be finitely piecewise affine,
1-Lipschitz, and satisfy
\[
 ax\le g(x)\le bx,\qquad 0<a<b\le1.
\]
An edge from \(n\) to \(m<n\) is admissible if \(2n-m\le N\).
Its cost is
\begin{equation}\label{humanprofile:edge}
 c_g(m,n)=g(m)-\min_{m\le x\le n}g(x).
\end{equation}
Costs are nonnegative. No intermediate endpoint is prescribed.

The two bounds we shall use are
\begin{align}
 C_1(a,b)&=\frac{b-a}{1+2b},
       \label{humanprofile:C1}\\
 C_2(a,b)&=\frac{(b-a)(2+a-2b-4ab)}
                    {(1-a)(1+2b)^2}.
       \label{humanprofile:C2}
\end{align}
Also put
\begin{equation}\label{humanprofile:parameters}
 b_c(a)=\frac{1+2a}{4+2a},\qquad
 R(a,b)=4(1+a)b^2-4(a+2)b+1+3a.
\end{equation}

\begin{theorem}[The profile estimate]\label{humanprofile:bound}
Suppose \(0<a\le1/4\), \(a<b\le1\), and let \(g\) satisfy the
preceding hypotheses. For every \(0\le n_0<N\), an admissible
chain from \(n_0\) to zero exists with total cost at most
\(NC(a,b)+\rho N\), for every prescribed \(\rho>0\),
where either of the following choices is valid:
\[
 \begin{array}{ll}
 C(a,b)=C_2(a,b),
 & b\le b_c(a)\ \hbox{and}\ R(a,b)\le0,\\[2pt]
 C(a,b)=C_1(a,b),
 & b\ge b_c(a).
 \end{array}
\]
The chain has at most
\begin{equation}\label{humanprofile:length}
 K\le C\left(1+\log\frac{N}{N-n_0}\right)
\end{equation}
edges, with an absolute constant. If a grid of mesh \(T\) contains
\(0,n_0,N\), its endpoints can be chosen on that grid with additional
cost \(O(KT)\). If the profile barriers hold with uniform additive
error \(\varepsilon N\), the additional cost is
\(O(K\varepsilon N+KT)\).
\end{theorem}

We first work on \([0,1]\). The arbitrary accuracy \(\rho\) avoids
any need to prove attainment of an infimum; the analytic argument has
a strict exponent margin and will choose \(\rho\) smaller than that
margin. The finite-chain and grid assertions are established at the end.

\subsection{Normalizing the endpoint and identifying the relevant gaps}

Replace \(g\) by
\[
 \widetilde g(x)=\min\{g(x),1+a-x\}.
\]
This is still 1-Lipschitz with the same two barriers, and
\(\widetilde g(1)=a\). Furthermore
\begin{equation}\label{humanprofile:normalization}
 c_g(m,n)\le c_{\widetilde g}(m,n).
\end{equation}
If \(g(m)\le1+a-m\), the left endpoint is unchanged and the
minimum can only decrease. Otherwise the new cost is at least
\((1+a-m)-(1+a-n)=n-m\), which bounds the old cost by
Lipschitz continuity. We may therefore assume \(g(1)=a\).

We use one published real-variable identity. In the notation below it
is precisely the hard-point formula of Keleti--Shmerkin
\cite[Lemma 5.4]{ks}:
if \(f\) is a piecewise affine 1-Lipschitz function with \(f(0)=0\),
then its infimum of total drops over partitions
\[
 1=t_0>t_1>\cdots\downarrow0,\qquad t_{i-1}\le2t_i,
\]
where a drop on \([t_i,t_{i-1}]\) uses its lower endpoint,
\(f(t_i)-\min_{[t_i,t_{i-1}]}f\), equals the sum of
\(f(u)-f(v)\) over the interval components \([u,v]\) of
\[
 \{t:f(t)=\min_{[t/2,t]}f\}.
\]
We now translate this identity carefully to the cost
\eqref{humanprofile:edge}.

Let \(\Phi_\zeta(g)\) be the infimum of chain costs from \(1-\zeta\)
to zero, and set \(f(t)=g(1-t)-a\). The edge \(n\to m\)
corresponds to \(x=1-n<y=1-m\), with \(y\le2x\).
Its cost is \(f(y)-\min_{[x,y]}f\). The corresponding
total-drop cost is \(f(x)-\min_{[x,y]}f\). Their difference is
\(f(y)-f(x)\), so telescoping gives
\[
 \Phi_\zeta(g)=T_\zeta(f)-g(1-\zeta),
\]
where \(T_\zeta(f)\) is the infimum of the usual total drop on
\([\zeta,1]\). Extending a finite partition below \(\zeta\) by
successive halvings adds at most \(\zeta\) to the cost.
Conversely, truncating an infinite partition at the interval crossing
\(\zeta\) changes that interval's cost by at most \(2\zeta\);
the ratio condition is preserved. Hence
\[
 T(f)-\zeta\le T_\zeta(f)\le T(f)+2\zeta,
 \qquad
 \Phi_0(g):=\lim_{\zeta\downarrow0}\Phi_\zeta(g)=T(f)-a.
\]

In the original coordinate the hard-point set is
\[
 H_g=\{x:g(x)=\min_{[x,(1+x)/2]}g\}.
\]
For a piecewise affine function this is a finite union of closed intervals,
possibly including degenerate intervals. Both zero and one belong to it.
The function \(g\) is nondecreasing on every component: hardness compares
sufficiently close pairs in that interval, and subdivision compares
arbitrary pairs. A component \([p,q]\) becomes
\([1-q,1-p]\) under reversal. The quoted identity therefore reads
\begin{equation}\label{humanprofile:hardidentity}
 \Phi_0(g)=
 \sum_{[p,q]\text{ component of }H_g}\bigl(g(q)-g(p)\bigr)-a.
\end{equation}
The subtraction of \(a\) is essential.

\begin{lemma}[The geometry of a gap]\label{humanprofile:gap}
If \((q,p)\) lies between consecutive components of \(H_g\), then
\[
 g(x)\ge g(p)\quad(q<x<p),\qquad g(q)\ge g(p).
\]
Its drop \(\delta=g(q)-g(p)\) is nonnegative; if it is positive,
\begin{equation}\label{humanprofile:gapbound}
 \delta\le p-\frac{1+q}{2}.
\end{equation}
\end{lemma}
\begin{proof}
Suppose \(g(x)<g(p)\) for some interior point. A minimum point \(z\)
of \(g\) on \([x,p]\) lies before \(p\) and has value below \(g(p)\).
On \([z,p]\) the profile is at least \(g(z)\). Any remaining part
of \([z,(1+z)/2]\) lies in \([p,(1+p)/2]\), where hardness of
\(p\) gives \(g\ge g(p)>g(z)\). Thus \(z\) would be hard,
contradicting the choice of the gap. Continuity proves the endpoint
inequality. If the drop is positive, hardness of \(q\) forces
\(p>(1+q)/2\) and \(g((1+q)/2)\ge g(q)\). Lipschitz
continuity from \((1+q)/2\) to \(p\) gives
\eqref{humanprofile:gapbound}.
\end{proof}

Telescoping the changes along all components and gaps, using
\(g(0)=0\), \(g(1)=a\), turns \eqref{humanprofile:hardidentity} into
\begin{equation}\label{humanprofile:drops}
 \Phi_0(g)=\sum_{\text{gaps}}\delta.
\end{equation}
We may discard gaps with zero drop.

\begin{lemma}[Controlling the remaining gaps]\label{humanprofile:tail}
If \(g(p)=v\), the sum \(S\) of drops of all positive gaps
contained in \([p,1]\) satisfies
\[
 S\le\frac{1-p}{2},\qquad
 S\le\frac{1-p+v-a}{2}.
\]
Consequently, for \(0\le\chi\le1/2\),
\begin{equation}\label{humanprofile:tailblend}
 S\le\frac{1-p}{2}+\chi(v-a).
\end{equation}
\end{lemma}
\begin{proof}
A gap \((q_i,p_i)\) has
\[
 \delta_i\le p_i-\frac{1+q_i}{2}
 \le\frac{p_i-q_i}{2}.
\]
The intervals are disjoint, proving the first estimate. Independently,
the drop is at most the negative variation of \(g\) on that gap.
On \([p,1]\), total variation is at most \(1-p\), and the
net change is \(a-v\). Its negative variation is therefore at most
\((1-p+v-a)/2\). A convex combination of the two estimates with
weights \(1-2\chi\) and \(2\chi\) proves the last statement.
\end{proof}

If there is only one positive gap, the barriers and
\eqref{humanprofile:gapbound} give
\[
 \delta\le p-\frac{1+q}{2},\qquad \delta\le bq-ap.
\]
Multiplying the first inequality by \(2b\), adding the second, and
dividing by \(1+2b\) yields
\begin{equation}\label{humanprofile:onegap}
 \delta\le\frac{(2b-a)p-b}{1+2b}\le C_1(a,b).
\end{equation}
With no positive gaps the cost is zero. We henceforth consider the first
two positive gaps. Write
\[
 (q_1,p_1),\ (q_2,p_2),\qquad
 u_i=g(q_i),\quad v_i=g(p_i),\quad \delta_i=u_i-v_i.
\]
The remaining gaps will be handled by Lemma~\ref{humanprofile:tail}.

\subsection{The bound when \(b\) is small}

Assume \(b\le b_c(a)\) and \(R(a,b)\le0\). In particular
\(b<1/2\). Put
\[
 s=1+2b,\qquad Z=(1-a)s^2,\qquad
 \chi=\frac{1+a-4ab+4b^2-2b}{Z}.
\]
The numerator is \(1-2b+a+4b(b-a)>0\), and
\[
 \chi-\frac12=\frac{R(a,b)}{2Z}\le0.
\]
Thus the blended tail estimate is legal.

The following seven inequalities follow respectively from the upper
barrier, lower barrier, gap bound, Lipschitz continuity, upper barrier,
gap bound, and the interval endpoint:
\[
 \begin{array}{rcl}
 u_1-bq_1&\le&0,\\
 ap_1-v_1&\le&0,\\
 q_1-2p_1+2\delta_1&\le&-1,\\
 u_2-v_1-q_2+p_1&\le&0,\\
 u_2-bq_2&\le&0,\\
 q_2-2p_2+2\delta_2&\le&-1,\\
 p_2&\le&1.
 \end{array}
\]
Multiply them, in that order, by
\[
 \begin{gathered}
 w_1=\frac1s,\quad
 w_2=\frac{1-2b}{(1-a)s},\quad
 w_3=\frac b s,\quad
 w_4=\frac{2b-a}{(1-a)s},\\
 w_5=\frac{1+2a-4b-2ab}{Z},\quad
 w_6=\frac{1-\chi}{2},\quad
 w_7=\frac12-\chi .
 \end{gathered}
\]
Every multiplier is nonnegative: the sign of \(w_5\) is exactly
\(b\le b_c(a)\), and the other signs follow from \(a<b<1/2\)
and \(0<\chi\le1/2\).
For transparency, the cancellation uses
\[
 \begin{gathered}
 w_3=bw_1,\quad w_1+2w_3=1,\quad
 aw_2-2w_3+w_4=0,\quad w_2+2w_3+w_4=1,\\
 w_6=w_4+bw_5,\quad w_4+w_5=\chi,\quad
 2w_6=1-\chi .
 \end{gathered}
\]
The resulting inequality is
\[
 \delta_1+\delta_2-\frac{p_2}{2}+\chi v_2
 \le-w_3-w_6+w_7=-\frac b{1+2b}-\frac\chi2.
\]
Apply \eqref{humanprofile:tailblend} after \(p_2\). Then
\[
 \begin{aligned}
 \Phi_0(g)
 &\le\delta_1+\delta_2+\frac{1-p_2}{2}+\chi(v_2-a)\\
 &\le\frac12-\frac b{1+2b}-\left(a+\frac12\right)\chi\\
 &=C_2(a,b).
 \end{aligned}
\]
The zero- and one-gap cases also satisfy this estimate because
\[
 C_2(a,b)-C_1(a,b)
 =\frac{(b-a)(1+2a-4b-2ab)}{(1-a)(1+2b)^2}\ge0.
\]

\subsection{The bound for \(b_c(a)\le b\le1/2\)}

Between \(p_1\) and \(q_2\), all hard components increase and
all omitted gaps have zero drop. Hence \(v_1\le u_2\).
Use the eight inequalities
\[
 \begin{array}{rcl}
 u_1-bq_1&\le&0,\\
 \delta_1-p_1+q_1/2&\le&-1/2,\\
 ap_1-v_1&\le&0,\\
 \delta_2-p_2+q_2/2&\le&-1/2,\\
 ap_2-v_2&\le&0,\\
 u_2-v_1-q_2+p_1&\le&0,\\
 p_2&\le&1,\\
 v_1-u_2&\le&0 .
 \end{array}
\]
Let \(Z_1=(1+2a)(1+2b)\). Their respective multipliers are
\[
 \begin{aligned}
 \omega_1&=\frac1{1+2b},&
 \omega_2&=\frac{2b}{1+2b},\\
 \omega_3&=\frac{2b}{Z_1},&
 \omega_4&=\frac{4b(1+a)}{Z_1},\\
 \omega_5&=\frac{1+2a-2b}{Z_1},&
 \omega_6&=\frac{2b(1+a)}{Z_1},\\
 \omega_7&=\frac{(6+8a)b-(1+2a)^2}{2Z_1},&
 \omega_8&=\frac{(2a+4)b-(2a+1)}{Z_1}.
 \end{aligned}
\]
All are nonnegative. The last sign is \(b\ge b_c(a)\).
The numerator of \(\omega_7\) increases with \(b\); at \(b=b_c(a)\)
it is
\[
 \frac{(1+a)(1-2a)(1+2a)}{a+2}>0.
\]
The remaining signs follow directly from \(0<a\le1/4\) and
\(b\le1/2\).

Expanding the weighted sum gives
\[
 \delta_1+\delta_2-\frac{p_2}{2}
 \le-\frac{\omega_2}{2}-\frac{\omega_4}{2}+\omega_7.
\]
Using the first tail estimate after \(p_2\), we obtain
\[
 \Phi_0(g)
 \le\frac12-\frac{\omega_2}{2}-\frac{\omega_4}{2}+\omega_7
 =C_1(a,b).
\]
The zero- and one-gap cases follow from \eqref{humanprofile:onegap}.

\subsection{The bound for \(b\ge1/2\)}

Here a direct potential is shorter than a gap calculation. Define
\[
 L(x)=2bx-b-g(x),\qquad
 V(x)=\frac{\max\{L(x),0\}}{1+2b}.
\]
The function \(L\) is nondecreasing, since \(g'\le1\le2b\)
almost everywhere. Moreover
\[
 (-g')_+\le\frac{2b-g'}{1+2b}.
\]
For \(g'\ge0\) the right side is nonnegative; for \(g'<0\),
the inequality reduces to \(g'\ge-1\).
Every edge cost is at most the negative variation of \(g\) on its
interval.

Above the zero region of \(L\), take maximal admissible downward
jumps, shortening the last one to reach that region. Their costs
sum to at most \(V(n_0)\). Once \(L(n)\le0\), if \(n>1/2\),
\[
 g(2n-1)\le b(2n-1)\le g(n).
\]
A minimum of \(g\) on \([2n-1,n]\), chosen leftmost, therefore
gives a zero-cost admissible continuation. When \(n\le1/2\), the
edge to zero is admissible and costs zero because \(g(0)=0\le g\).

To make termination precise, perform this construction on a uniform grid
and a piecewise affine interpolation. In the positive-\(L\) region stop
at the first grid point above its zero region and cross to the next
point below, at cost at most the mesh. The zero-cost minima then occur
at grid points and strictly decrease the endpoint, so the procedure
terminates. If the prescribed start is not on this grid, an initial
step to the nearest point below it costs at most the mesh and is
admissible once the mesh is small. Thus the grid construction costs
at most \(V(n_0)+O(\text{mesh})\). Compress the chain by the
procedure below. Its length is then bounded independently of the mesh,
so transferring it from the interpolant to \(g\) costs only
\(O(K\,\text{mesh})\). Taking the mesh sufficiently small makes
the total error less than the prescribed \(\rho\). Since
\[
 V(n_0)\le V(1)=\frac{b-a}{1+2b}=C_1(a,b),
\]
this proves the required high-\(b\) estimate.

\subsection{Finite chains and grid errors}

First, an infimum for a finite start is bounded by the limiting cost
used above. Truncating a chain from a larger starting point at its
first crossing of a smaller start does not increase its cost:
the crossing edge is replaced by a subinterval with the same lower
endpoint, whose minimum can only increase. Conversely, maximal
admissible edges bridge the two starting points at total cost at
most their distance, by the variation bound. Hence
\[
 0\le\Phi_0(g)-\Phi_\zeta(g)\le\zeta .
\]

Every finite chain can be compressed without increasing its cost.
For \(l<m<n\),
\[
 c_g(l,n)\le c_g(l,m)+c_g(m,n).
\]
Indeed, if \(A=\min_{[l,m]}g\) and \(B=\min_{[m,n]}g\), the
difference between the right side and left side is
\(g(m)-\max(A,B)\ge0\). Merge neighboring edges whenever their
union remains admissible. In the resulting chain, every pair of
successive edges more than doubles the complementary distance
from the endpoint one:
\[
 2n_i-n_{i+2}>1
 \quad\Longrightarrow\quad
 1-n_{i+2}>2(1-n_i).
\]
A start at \(1-\zeta\) therefore needs at most
\(2\lceil\log_2(1/\zeta)\rceil+O(1)\) edges.

Choose a finite chain within \(\rho\) of its finite-start infimum
and compress it. The resulting cost is at most \(C(a,b)+\rho\)
and its length has the displayed bound. No limiting chain or attainment
is needed. The high-\(b\) construction gives the same conclusion:
its interpolant differs uniformly from \(g\) by at most the mesh,
so each retained edge cost differs by at most twice the mesh.
Rescaling gives \eqref{humanprofile:length} and the asserted
cost \(NC(a,b)+\rho N\).

Finally round all endpoints downward to a grid of mesh \(T\) containing
\(0,n_0,N\). Admissibility is preserved because
\[
 \left\lfloor\frac mT\right\rfloor
 \ge\left\lfloor\frac{2n-N}{T}\right\rfloor
 \ge2\left\lfloor\frac nT\right\rfloor-\frac NT .
\]
Each endpoint moves at most \(T\), so each cost changes by at most
\(2T\). If the original barriers have additive error
\(\varepsilon N\), first clamp the profile between \(ax\) and \(bx\).
The result is still 1-Lipschitz and differs by at most
\(\varepsilon N\); a cost changes by at most twice that amount.
Endpoint normalization only increases the cost available to the original
profile. Summing these errors proves the final assertions of
Theorem~\ref{humanprofile:bound}.

\subsection{Evaluating the dimension cutoff}

Set
\[
 \alpha=\frac{3-\sqrt7}{4},\qquad
 a_c=\frac{\sqrt6-2}{4},\qquad
 b_F(a)=\frac{2a}{1-2a}.
\]
For \(\alpha\le a\le a_c\), define
\begin{equation}\label{humanprofile:root}
 r_2(a)=
 \frac{8a^2-a+2-
 \sqrt{64a^4+32a^3-63a^2-28a+4}}
 {4(1+4a-2a^2)}.
\end{equation}
The radical is real on this interval, as verified below. Rationalizing the
numerator gives the expression in \eqref{human:root}: if
\(A=2-a+8a^2\) and \(\Delta=4-28a-63a^2+32a^3+64a^4\), then
\(A^2-\Delta=24a(1+4a-2a^2)\).
The final curve is
\begin{equation}\label{humanprofile:curve}
 B_{\mathrm H}(d)=
 \begin{cases}
  2d-1,&1<d\le1+\alpha,\\[2pt]
  1+r_2(d-1),&1+\alpha<d\le1+a_c,\\[2pt]
  \dfrac1{3-2d},&1+a_c<d\le5/4.
 \end{cases}
\end{equation}

\begin{proposition}[What the inequality supplies]
\label{humanprofile:algebra}
Let \(1<d\le5/4\) and \(d\le D<B_{\mathrm H}(d)\).
Then either
\[
 D<2d-1,
\]
or there are \(s<d\), \(u>D\), with \(1<s<5/4\),
\(s<u<2\), for which one of the two parameter conditions in
Theorem~\ref{humanprofile:bound} holds and
\(C(s-1,u-1)<s-1\).
\end{proposition}
\begin{proof}
Put \(a=d-1\), \(b=D-1\). Only \(b\ge2a\) needs examination.
First let
\(\alpha\le a\le a_c\) and \(2a\le b\le b_c(a)\).
The guard \(R(a,b)<0\) holds throughout this region. Indeed
\[
 \partial_bR=8(1+a)b-4(a+2)\le-4\qquad(b\le1/2),
\]
and
\[
 R(a,2a)=16a^3+8a^2-13a+1.
\]
Its derivative is \(48a^2+16a-13\le-6\) on \([0,1/4]\),
whereas
\[
 R(\alpha,2\alpha)=3(11\alpha-1)<0.
\]
For the last sign, \(\alpha<1/11\) is equivalent to
\(29<11\sqrt7\), and squaring gives \(841<847\).

The derivative of \(C_2\) in \(b\) is
\[
 \partial_bC_2=
 \frac{2+11a+8a^2-(8+14a+8a^2)b}
      {(1-a)(1+2b)^3}>0
 \qquad(b\le b_c(a)).
\]
The numerator decreases with \(b\); its value at \(b_c(a)\),
multiplied by \(4+2a\), is \(18a(1+a)>0\).
At the lower endpoint,
\[
 C_2(a,2a)-a
 =\frac{a(1+2a)(8a^2-12a+1)}
        {(1-a)(1+4a)^2}.
\]
This is zero at \(a=\alpha\) and negative for
\(\alpha<a\le a_c\).
At \(b=b_c(a)\), the formulas \(C_1,C_2\) agree. Also
\[
 C_1(a,b)<a
 \quad\Longleftrightarrow\quad b<b_F(a).
\]
The equality \(b_F(a)=b_c(a)\) is equivalent to
\(8a^2+8a-1=0\), hence to \(a=a_c\).
For \(a<a_c\), \(b_F(a)<b_c(a)\), so
\(C_2(a,b_c(a))>a\).

There is therefore one root of \(C_2(a,b)=a\) in
\([2a,b_c(a)]\), with the indicated endpoint equalities.
Clearing its positive denominator gives
\[
 (4a^2-8a-2)b^2+(8a^2-a+2)b-3a=0.
\]
The smaller root is \eqref{humanprofile:root}. The crossing has
strictly positive derivative, so its discriminant is positive.
This proves the reality and continuity of the radical expression.
Its endpoint values are
\[
 r_2(\alpha)=2\alpha,\qquad
 r_2(a_c)=b_c(a_c)=\frac{\sqrt6-1}{5}.
\]
Thus the middle branch gives exactly the required strict cost
inequality above the coherent range, and the curve is continuous
at both joins.

Now let \(a_c\le a\le1/4\).
For \(b\ge b_c(a)\), Theorem~\ref{humanprofile:bound} gives
\(C_1<a\) exactly when \(b<b_F(a)\).
If \(b<b_c(a)\), increase the upper profile barrier to
\(b_c(a)x\). For \(a>a_c\), the preceding crossing calculation
gives \(b_c(a)<b_F(a)\), so the high-side estimate at that
larger barrier has
\[
 C_1(a,b_c(a))<a.
\]
At \(a=a_c\) and \(b<b_c(a)\), the strict low-side estimate
has already been proved in the middle branch.
Finally \(1+b_F(a)=1/(1-2a)=1/(3-2d)\), proving the last
branch of \eqref{humanprofile:curve}.

All inequalities used by the profile branch are strict in cost.
The guard is also strict wherever the low branch is used here.
Away from the common boundary, continuity permits a small decrease
of \(a\) and increase of the chosen upper barrier. At the common
boundary one may move to its high side, since \(C_1=C_2\) there.
In the enlarged-barrier case \(a>a_c,\ b<b_c(a)\), choose
\(a'<a\) sufficiently close that \(a'>a_c\) and
\(b_c(a')>b\), and use the upper barrier \(b_c(a')\).
These choices give \(s=1+a'<d\) and some \(u>D\) with the
required strict profile bound. They may be chosen with
\(1<s<5/4\), \(s<u<2\), proving the assertion.
\end{proof}

The analytic proof can now use the coherent implication when
\(D<2d-1\). In every remaining case it uses
Theorem~\ref{humanprofile:bound} with the strict parameter slack
provided by Proposition~\ref{humanprofile:algebra}. This is the
complete profile input needed for the curve
\eqref{humanprofile:curve}.
